% Technical report, course project 5.5: Attacking LLMs via subtle paraphrasing.
% ACM double-column conference format (acmart, sigconf). 4-8 pages.
% Compile on Overleaf (acmart is built in) or locally with: latexmk -pdf main.tex
\documentclass[sigconf,nonacm]{acmart}

\usepackage{booktabs}

\settopmatter{printacmref=false}
\renewcommand\footnotetextcopyrightpermission[1]{}
\pagestyle{plain}

\begin{document}

\title{Attacking LLMs via Subtle Paraphrasing}

% TODO: names and emails of all three team members
\author{Guy Zilberman}
\affiliation{\institution{Tel Aviv University}\country{Israel}}
\author{Team member 2}
\affiliation{\institution{Tel Aviv University}\country{Israel}}
\author{Team member 3}
\affiliation{\institution{Tel Aviv University}\country{Israel}}

\begin{abstract}
% ~200 words or fewer. Problem, method, main numbers, takeaway.
% Main numbers to use: 96 curated NQ-Open questions; genuine confirmed wrong
% answers on 22 of 95 attacked questions; reasoning vs non-reasoning victim
% (0 vs 21 genuine on 16 shared questions); B1 plain rewording vs full method.
TODO
\end{abstract}

\maketitle

\section{Introduction}
% Motivation from the project brief: LLMs are used as an authority, e.g. for
% fact checking. Question: does a semantically equivalent phrasing exist that
% leads to a factually wrong answer with high probability (no typos, no
% invisible characters)? What we did, at a high level, and the main findings.
TODO

\section{Related Work}
% Prior work on prompt / paraphrase sensitivity of LLMs and adversarial text
% attacks; GEPA \cite{agrawal2025gepa} for reflective prompt evolution, which
% our search builds on. Where our work fits in.
TODO

\section{Technical Approach}
% Pipeline (figure): attacker -> free checks (leaked notes, >=60% word
% overlap, answer leak) -> equivalence and answer-preservation checks ->
% victim (1 greedy + 4 sampled answers) -> LLM judge -> fitness.
% Search: 8 rounds x 5 paraphrases; evolutionary parent selection (one parent
% per distinct wrong answer, after GEPA's Pareto idea) vs reflective.
% Minimal-edit attacker prompt. Confidence groups (20 samples on the
% original) and the 20-answer re-test with a one-sided Fisher exact test.
% Models and why: victim qwen3:4b-instruct / qwen3:4b, attacker llama3.1:8b,
% judges gemma3:12b (judge accuracy from tests/eval_judges.py).
% Question curation (data/CURATION.md).
TODO

\section{Results}
\subsection{Experimental setup}
% Datasets (22 / 16 / 96 questions), models, hardware (A100, Ollama 0.34.4),
% metrics: raw success, robust, confirmed, genuine (manual review), and
% the second-annotator agreement.
TODO

\subsection{Main study}
% Table: full method vs B1 on the 96 questions: raw, confirmed, genuine,
% questions with a genuine success; by confidence group.
TODO

\subsection{What each part of the method adds}
% Table: B1-B4 on the 16-question set.
TODO

\subsection{Reasoning vs non-reasoning victim}
% 0 vs 21 genuine successes on the same 16 questions; matching wrong answers.
TODO

\section{Discussion}
% Drift and reward hacking (what the search exploits), judge errors,
% confirmed != genuine, hardware nondeterminism (GPU / Ollama version),
% small question sets, single attacker and victim family. Future work.
TODO

\section{Conclusion}
TODO

\section{Contributions}
% What each team member did.
TODO

\bibliographystyle{ACM-Reference-Format}
\bibliography{references}

\end{document}
